% Supplementary Information of the ERL Letter (D-107): Methods of the near-term emergence
% analysis (H3), new for the Supplementary. Design: PREREG section 10, D-081, D-090, D-095
% (ATL24, Amendment 17), D-096 (sensitivity forms). Every number is a \metric{} macro.

\subsection*{Near-term emergence and exposure}
\emph{Reference level and scenarios.} The 2026 reference level is the mean of the \metric{h3.level_2026.n_records} DAHITI records of 2026, \metric{h3.level_2026_m}~m on the DAHITI reference (\metric{erl.text.h3.level_2026_bhs_m}~m in BHS-77); the latest three-record running median is a sensitivity. The scenarios are declines of \metric{h3.scenarios_text}~m below it; the pre-registered decision rule uses the \metric{h3.n_near_term} near-term declines of 0.5, 1 and 2~m.

\emph{Surfaces.} The \emph{uncorrected} surface is each grid as the hindcast reads it on the DAHITI reference through the per-cell datum chain (the total variant), with no constant and no error model: the grid as users apply it. The \emph{corrected} surface applies the registered error model. Its input is, per grid, observed unit (epoch and level bin) and sector, the vertical shift that minimises the absolute total area bias, searched on the grid of the hindcast's shape constant with the same tie rules; it is modelled as a function of depth $d = L_{\min} - z$ measured on the grid, where $L_{\min}$ is the lowest fitted level. The registered form, \texttt{linear\_hold}, is a Theil--Sen line in $d$ fitted on E2 and E3 and held constant beyond the fitted depth range; a slope with $1 + b < 0.2$ falls back to the constant form, and a sector with fewer than five usable units uses the domain model. The registered sensitivities are the \texttt{constant} form (each sector's median fitted shift, inside the fitted shifts by construction) and the \texttt{linear} form. The \emph{fused} surface is the waterline elevation model where it is finite and the cell is not known sea at the lowest observed unit, and the corrected grid elsewhere, flagged as extrapolation; where the grid contradicts the lowest observed extent it is clamped at $L_{\min}$, and such cells are counted.

\emph{Emerged area.} For a decline $\Delta$, $E(\Delta) = A(L_{2026}) - A(L_{2026} - \Delta)$, where $A(L)$ is the area of the surface's flood fill from the open sea at level $L$ inside the domain. Each surface is thus referenced to its own 2026 shoreline; the observed 2026 extent as a common baseline is a sensitivity.

\emph{Intervals.} Split-conformal 90\% intervals\citep{lei2018} are calibrated on the absolute errors of the emerged area at held-out observed units of the leave-one-epoch-out hindcast (fit on two epochs, predict the third), at observed levels only and without each epoch's top unit: \metric{h3.n_scores.gebco2026.uncorrected} scores per grid and surface. They are calibrated per grid, surface, region and quantity, never summed, clipped to zero and to the baseline area, and infinite with fewer than nine scores. The corrected and fused surfaces use the corrected grid's scores. The interval from the scores of the lowest held-out epoch alone is a sensitivity.

\emph{Transfer and decision rule.} For each held-out epoch, the transfer test compares the median absolute emergence error of the corrected and the uncorrected grid at that epoch's scored units, and reports the median error of the transferred shift. The correction is decision-relevant at a decline if the corrected estimate lies outside the 90\% interval of the uncorrected one, or the reverse. Per grid and for the domain, H3 is ``supported'' if the correction is decision-relevant at all three near-term declines and ``partial'' if at some; grids are never pooled, and exposure does not enter the rule.

\emph{Exposure.} From OpenStreetMap\citep{osm}: oil and gas platforms and wells newly onshore (\metric{h3.exposure.platforms.n_features} features in the domain); pipeline (\metric{h3.exposure.pipelines.n_features}) and channel (\metric{h3.exposure.channels.n_features}) length that emerges; ports, piers and ferry terminals (\metric{h3.exposure.ports.n_features}) that have the sea within \metric{h3.port_access_km}~km in 2026 but not at the scenario level; and settlements (\metric{h3.exposure.settlements.n_features}) with the sea within \metric{h3.coast_km}~km in 2026 but not at the scenario level. Habitat stand-ins with open licences: the herbaceous wetland class of ESA WorldCover 2021\citep{worldcover2021}, the wetland classes of GWL\_FCS30 2020 and the wetland classes of the Global Lakes and Wetlands Database version~2\citep{zhang2023,lehner2025}. Exposure is counted at the scenario level and at the levels equivalent to the bounds of the area interval, rounded outward. Protected areas, Ramsar sites, Key Biodiversity Areas, Important Marine Mammal Areas and the range of the Caspian seal were not used in this analysis: their terms require acceptance or registration by the PI, and OpenStreetMap holds too few protected areas in the domain.

\emph{Scenario dates.} The primary pathways are the best-fit SSP2-4.5 and SSP5-8.5 totals and the observed 1996--2021 trend of \citet{samant2023}, and the medians of the low, medium, medium-to-high and high emission scenarios of \citet{hoseini2025}. Each pathway's own change after its base year is applied from the 2026 level, linearly between the years its source states and never beyond 2100. The quotes behind every pathway were checked against the retrieved sources. Range and sensitivity pathways are tabulated beside them.

\emph{Independent seabed check (pre-registration Amendment~17).} ICESat-2 ATL24 nearshore bathymetry (version~2)\citep{atl24_v2,is2_atl24_atbd_v2} was to test the surfaces below the 2026 waterline. The \metric{h3.atl24.n_granules_found} granules that the NASA Common Metadata Repository lists for the domain's bounding box (\metric{h3.atl24.size_mb_cmr}~MB) were read in place: the data extent of \metric{h3.atl24.n_granules_extent_in_box} of them meets the box, all extents lying between \metric{h3.atl24.extent.lon_min_deg}\textdegree{} and \metric{h3.atl24.extent.lon_max_deg}\textdegree{}E and \metric{h3.atl24.extent.lat_min_deg}\textdegree{} and \metric{h3.atl24.extent.lat_max_deg}\textdegree{}N (the Persian Gulf), and \metric{h3.atl24.spot_check.n_with_box_photons} of \metric{h3.atl24.spot_check.n_sampled} granules drawn at random for a full read hold a photon in the box. Every verdict of the check is therefore ``\metric{h3.atl24.gebco2026.verdict}''. The documents read do not say why the Caspian is absent.
